\documentclass[twoside]{article}
\usepackage{package}
\usepackage[T1]{fontenc}
\usepackage{subfiles}
% Silenciar la advertencia de versalita + negrita inexistente
\WarningFilter{latex}{Some font shapes were not available}
\WarningFilter{latex}{Font shape `T1/lmr/bx/sc'}

% Silenciar la advertencia de caracteres perdidos en el nullfont (para 01_introduccion.tex)
\WarningFilter{latex}{Missing character: There is no}
%Nuevos comandos
\renewcommand{\thefootnote}{\arabic{footnote}}
\raggedbottom



\begin{document}
\begin{center}
    \thispagestyle{empty}
    \large{\textsc{Personalizar Matlab y Sistemas Lineales invariantes en el Tiempo}}\\[0.5cm] 
    \large{Fabrizio Benavidez\footnote[1]{\texttt{fabrizio.benavidez@mi.unc.edu.ar}}}\\[0.5cm]

    \large{\footnotemark[1]\textit{Facultad de Ciencias Exactas, Físicas y Naturales, Ingeniería Aeronáutica}}
    \\[0.5cm]

    \begin{minipage}{0.7\textwidth}
    \hrule
    \vspace{0.2cm}
        \normalsize{\textsc{Resumen}En el siguiente trabajo se conoce con más profundidad la interfaz de MATLAB y algunas facilidades. Además, se trabaja con sistemas lineales: ceros, polos y ganancia, así como su representación en un espacio de estados. Se analiza distintas funciones de transferencias y sus diagramas de Bode, Nyquist y gráficas de respuestas entradas e impulsos.}\\[0.2cm]
        \normalsize{\textit{Palabras claves}: sistemas, MATLAB, funciones, Bode, Nyquist\ }
    \vspace{0.2cm}
    \hrule
    \end{minipage}
    
    \vfil

    \includegraphics[width = 0.4\textwidth]{imagenes/Matlab_Logo.png}\\[0.5cm]

    \vfil

    \today\\[0.2cm]
    \large{Semana 8}
\end{center}







\newpage

\section*{Guía 16: Sistemas Lineales invariantes en el Tiempo} 
\subfile{secciones/guia16}














\newpage
% \thispagestyle{empty}
% \bibliography{biblioteca.bib}
% \bibliographystyle{plain}


\end{document}